% Standalone article section for the parent manuscript.
% Requires amsmath, amssymb, amsthm, hyperref and ordinary theorem/corollary/remark.
% New macros: \gsZ, \gsJ, \gsB, \gsV, \gsQ, \gsE, \gsT, \gsL.
% Bibliography keys: ShpotParis, Connon, DLMF, RepoTailKernels,
% RepoBilinear.  Proposed bibliography text is supplied separately.
\newcommand{\gsZ}{\mathcal Z}
\newcommand{\gsJ}{\mathcal J}
\newcommand{\gsB}{\mathcal B}
\newcommand{\gsV}{\mathcal V}
\newcommand{\gsQ}{\mathcal Q}
\newcommand{\gsE}{\mathcal E}
\newcommand{\gsT}{\mathcal T}
\newcommand{\gsL}{\mathcal L}

\section{Independent nonperiodic shifts: a completed product generator}
\label{gs:section}

The ordinary integral
\[
 \int_0^1\gamma_m^{(r)}(a+x)\gamma_n^{(t)}(b+x)\,dx,
 \qquad \Re a,\Re b>0,
\]
has two independent argument shifts.  It is a different object from a
periodically translated product.  Question R8 of \emph{Tail Kernels and
Boundary Products} asks for an order-generating formula, especially its
antisymmetric first-argument-derivative component
\cite{RepoTailKernels}.  This section supplies such a formula in an
explicit function class: finitely many Gauss hypergeometric terms,
holomorphic Gamma divided differences, and two geometrically convergent
single Hurwitz-zeta series.  It also gives a smaller function class for
all products of ordinary polygamma functions.

The mechanism extends the classical product-integral method of
Shpot and Paris \cite{ShpotParis}.  Their coincident shift is the
case $a=b=1$ before the spectral poles are removed.  The logarithmic
series occurring in digamma-square integrals is also classical; see
Connon \cite{Connon}.  The additions established here are the
independent-shift completion as a \emph{joint} holomorphic germ, its
explicit antisymmetric derivative, and the finite closure of all
ordinary argument derivatives in one logarithmic Hurwitz family.
These are claims relative to the inspected project sources; historical
priority for every specialization is not asserted.

\subsection{Unfolding the two-dimensional lattice}

For $\Re A,\Re B>0$ and $\Re(s+t)>1$, define the elementary
incomplete-beta kernel
\begin{align}
 \gsB_{s,t}(A,B)
 &=\int_0^\infty\frac{dy}{(y+A)^s(y+B)^t}\notag\\
 &=\frac{B^{1-s-t}}{s+t-1}\,
 {}_2F_1\left(s,s+t-1;s+t;1-\frac AB\right).
 \label{gs:beta-kernel}
\end{align}
The first expression fixes the branches of the second.  Thus the
formula is regular when $A=B$, with value
$A^{1-s-t}/(s+t-1)$.

\begin{theorem}[Exact diagonal unfolding]\label{gs:unfolding}
For $\Re s,\Re t>1$ and $\Re a,\Re b>0$,
\begin{equation}\label{gs:unfolded}
 \int_0^1\zeta(s,a+x)\zeta(t,b+x)\,dx
 =\sum_{k=0}^\infty\gsB_{s,t}(a+k,b)
  +\sum_{k=1}^\infty\gsB_{s,t}(a,b+k).
\end{equation}
Both sums converge absolutely and locally normally in this domain.
\end{theorem}

\begin{proof}
Expand both Hurwitz functions by their absolutely convergent Dirichlet
series.  On the lattice half-plane $m\ge n\ge0$, put $m=n+k$.
For each fixed $k\ge0$, the intervals $n+x$, $0\le x\le1$, tile the
positive real axis; consequently
\[
 \sum_{n=0}^\infty\int_0^1
 (x+n+k+a)^{-s}(x+n+b)^{-t}\,dx
 =\gsB_{s,t}(a+k,b).
\]
The complementary half-plane $n>m\ge0$ gives the second sum, starting
at $k=1$.  This starting index is essential: the lattice diagonal is
counted once.  Absolute convergence follows by replacing the complex
shifts with lower bounds for their real parts and applying the same
argument to positive integrands.  The bounds are uniform on compact
subsets, proving normal convergence.
\end{proof}

Choose integers $K\ge0$ and $L\ge1$, and write
\begin{equation}\label{gs:cd}
 c=K+a-b,\qquad d=L+b-a.
\end{equation}
They may always be chosen so large that
\begin{equation}\label{gs:cutoff-condition}
 \Re c>|b|,\qquad \Re d>|a|.
\end{equation}
The choices can be kept fixed on any sufficiently small neighborhood
of given shifts.  Put
\begin{equation}\label{gs:finite-prefix}
 \Phi_{s,t}(a,b)=
 \sum_{k=0}^{K-1}\gsB_{s,t}(a+k,b)
 +\sum_{k=1}^{L-1}\gsB_{s,t}(a,b+k),
\end{equation}
where an empty sum is zero.  A useful accelerated form of
\eqref{gs:unfolded} is
\begin{align}
 I(s,t;a,b)={}&\Phi_{s,t}(a,b)
 +B(1-t,s+t-1)\zeta(s+t-1,c)\notag\\
 &+B(1-s,s+t-1)\zeta(s+t-1,d)\notag\\
 &-\sum_{j=0}^\infty
 \frac{(-1)^j(s)_j b^{j+1-t}}{j!(j+1-t)}\zeta(s+j,c)\notag\\
 &-\sum_{j=0}^\infty
 \frac{(-1)^j(t)_j a^{j+1-s}}{j!(j+1-s)}\zeta(t+j,d),
 \label{gs:meromorphic-master}
\end{align}
where $I$ denotes the integral on the left of \eqref{gs:unfolded}.
Here $B$ without calligraphic type is Euler's beta function.
Identity \eqref{gs:meromorphic-master} is understood meromorphically;
its separate terms need not be regular at integral orders.

Indeed, for a fixed $q$ with $\Re q>|b|$, the change of variable
$y+b=z$ in \eqref{gs:beta-kernel}, followed by subtraction of the
interval $0<z<b$, gives
\begin{equation}\label{gs:strip-evaluation}
 \gsB_{s,t}(q+b,b)=q^{1-s-t}B(1-t,s+t-1)
 -\sum_{j=0}^\infty
 \frac{(-1)^j(s)_j b^{j+1-t}}{j!(j+1-t)}q^{-s-j}.
\end{equation}
This first holds where the beta integral and the subtracted interval
both converge, and then holds by meromorphic continuation.  Summing
over $q=c,c+1,\ldots$ proves \eqref{gs:meromorphic-master}; the other
tail follows by exchange of the two factors.  Alternatively one may
first subtract the $j=0$ term and continue the resulting normally
convergent remainder.  In either derivation, the condition
\eqref{gs:cutoff-condition} gives a geometric majorant for the
$j$-series on compact sets avoiding the displayed poles.

\subsection{Removal of both spectral poles without an iterated limit}

Use the standard convention
\begin{equation}\label{gs:stieltjes-convention}
 \gsZ(u,z)=\zeta(1+u,z)-\frac1u
 =\sum_{m=0}^\infty\frac{(-1)^m\gamma_m(z)}{m!}u^m.
\end{equation}
For $\Re z>0$ this is entire in $u$.  The required ordinary product
generator is
\begin{equation}\label{gs:target-master}
 \gsJ(u,v;a,b)=\int_0^1\gsZ(u,a+x)\gsZ(v,b+x)\,dx.
\end{equation}
We now evaluate it; its definition alone is not the claimed result.
The Hurwitz recurrence and argument-derivative formulas used below
have the conventions of DLMF, \S25.11 \cite{DLMF}.

Define the universal Gamma functions
\[
 \gsQ(u,v)=\frac{\Gamma(1+u+v)}{\Gamma(1+u)\Gamma(1+v)},
 \qquad \gsE(w)=\frac{\pi w}{\sin\pi w}.
\]
The following quotient has a joint removable singularity:
\begin{equation}\label{gs:universal-quotient}
 \gsT(u,v)=
 \frac{u+v-\gsQ(u,v)\{u\gsE(v)+v\gsE(u)\}}
 {uv(u+v)}.
\end{equation}
For an explicitly holomorphic interpretation, write
\begin{align}
 \gsT(u,v)&=-\frac{\gsQ(u,v)-1}{uv}
 -\gsQ(u,v)\sum_{h=1}^\infty e_{2h}
       \frac{u^{2h-1}+v^{2h-1}}{u+v},\label{gs:universal-holomorphic}\\
 e_{2h}&=2(1-2^{1-2h})\zeta(2h).
\end{align}
Each quotient in the series is a polynomial; also
$\gsQ-1$ is divisible by $uv$ in the ring of convergent power
series.  In particular $\gsT(0,0)=-2\zeta(2)$.

The only infinite series in the completed answer are
\begin{equation}\label{gs:V}
 \gsV(u,v;c,z)=
 -\sum_{j=1}^\infty
 \frac{(-1)^j(1+u)_j z^{j-v}}{j!(j-v)}
                    \zeta(1+u+j,c),\qquad \Re c>|z|.
\end{equation}
They are holomorphic near $(u,v)=(0,0)$, locally normally in all
variables.  No term at $j=0$ is included.

\begin{theorem}[Completed independent-shift generator]
\label{gs:completed}
Under \eqref{gs:cd}--\eqref{gs:cutoff-condition}, the ordinary integral
\eqref{gs:target-master} is the following joint holomorphic germ:
\begin{align}
 \boxed{\begin{aligned}
 \gsJ(u,v;a,b)={}&\Phi_{1+u,1+v}(a,b)+\gsT(u,v)\\
 &+\frac{b^{-v}\gsZ(u,c)
       -\gsQ(u,v)\gsE(v)\gsZ(u+v,c)}{v}\\
 &+\frac{a^{-u}\gsZ(v,d)
       -\gsQ(u,v)\gsE(u)\gsZ(u+v,d)}{u}\\
 &+\gsV(u,v;c,b)+\gsV(v,u;d,a).
 \end{aligned}}
 \label{gs:completed-formula}
\end{align}
Every displayed divided difference has its unique holomorphic value
on its zero-denominator face.  All mixed spectral and argument
derivatives may be taken term by term in the two $\gsV$ series.
The result is independent of the permissible integers $K,L$.
\end{theorem}

\begin{proof}
The unit-interval integral of one Hurwitz function is
\[
 \int_0^1\zeta(s,a+x)\,dx=\frac{a^{1-s}}{s-1}.
\]
Consequently, with $s=1+u$, $t=1+v$,
\begin{equation}\label{gs:pole-removal}
 \gsJ(u,v;a,b)=I(1+u,1+v;a,b)
             -\frac{a^{-u}+b^{-v}-1}{uv}.
\end{equation}
In \eqref{gs:meromorphic-master}, isolate $j=0$ in each series and
use
\[
 B(-v,1+u+v)=-\frac{\gsQ(u,v)\gsE(v)}v,
 \qquad
 B(-u,1+u+v)=-\frac{\gsQ(u,v)\gsE(u)}u.
\]
Replacing each Hurwitz function with its pole plus \eqref{gs:stieltjes-convention}
leaves precisely the two divided differences in
\eqref{gs:completed-formula}.  The remaining rational pole terms
combine to \eqref{gs:universal-quotient}.  Formula
\eqref{gs:universal-holomorphic} proves that this last apparent
singularity is jointly removable, including on $u+v=0$.
The numerator of the first divided difference vanishes identically
on $v=0$; that of the second vanishes on $u=0$.

For normal convergence, on a compact subset of
$\Re c>|z|$ choose $\rho<1$ with $|z|/\Re c\le\rho$.
For large $j$, the defining Hurwitz series gives a bound of the form
$C(\Re c)^{-j}$ times a fixed polynomial in $j$; the Pochhammer
ratio $(1+u)_j/j!$ also has polynomial growth uniformly for $u$ in
a compact neighborhood of zero.  Thus the summands in
\eqref{gs:V} are bounded by $C'j^M\rho^j$.  Cauchy's estimates
on a slightly larger compact set give the same conclusion after any
fixed mixed differentiation.  The finite prefix is holomorphic
because $\Re(2+u+v)>1$ near the origin.
The meromorphic identity therefore extends to a holomorphic equality.
Its left side is independent of $K,L$, proving the final assertion.
\end{proof}

\begin{corollary}[Every independent-shift Stieltjes product]
\label{gs:all-indices}
For all $m,n,r,t\ge0$ and $\Re a,\Re b>0$,
\begin{equation}\label{gs:all-indices-formula}
 \int_0^1\gamma_m^{(r)}(a+x)\gamma_n^{(t)}(b+x)\,dx
 =(-1)^{m+n}
 \left.\partial_a^r\partial_b^t\partial_u^m\partial_v^n
                  \gsJ(u,v;a,b)\right|_{u=v=0},
\end{equation}
where the right side is evaluated by the explicit convergent formula
\eqref{gs:completed-formula}.  All constants of integration and all
spectral regularizations are fixed by this formula.
\end{corollary}

This is a convergent reduction, not a claim of finite reduction to
ordinary zeta values at arbitrary shifts.  To obtain a prescribed
coefficient, expand the finite Gamma germ through the required degree,
use generalized Stieltjes constants at $c,d$, and extract that degree
in each geometrically convergent series.  The resulting coefficients
involve finite harmonic sums from $(1+u)_j$, powers of $\log a$ and
$\log b$, and Hurwitz order derivatives.  Each requested coefficient
therefore has an explicit, rapidly convergent single-sum expression.
More explicitly, the entire generalized-harmonic contribution is
encoded by the finite-product identity
\[
 \frac{(1+u)_j}{j!}
 =\exp\left\{\sum_{k=1}^\infty
        \frac{(-1)^{k+1}}{k}H_j^{(k)}u^k\right\},
 \qquad H_j^{(k)}=\sum_{h=1}^j h^{-k},\quad |u|<1.
\]

\subsection{The finite terms are polylogarithmic at integral orders}

There is no need to introduce parameter derivatives of a hypergeometric
function as additional undefined constants.  For positive $A>B$, put
$q=A-B$ and $z=q/A$.  Equation \eqref{gs:beta-kernel} becomes
\begin{equation}\label{gs:polylog-kernel}
 \gsB_{1+u,1+v}(A,B)
 =q^{-1-u-v}\int_0^z
       \frac{\exp\{(u+v)\log y-v\log(1-y)\}}{1-y}\,dy.
\end{equation}
Every fixed Taylor coefficient is a finite linear combination of
harmonic polylogarithms on the alphabet $\{0,1\}$ and logarithms;
this follows by the shuffle product for powers of $\log y$ and
$\log(1-y)$ followed by one integration with kernel $1/(1-y)$.
The coefficient of total degree $N$ has weight at most $N+1$ before
multiplication by powers of $\log q$.  Other positive $A,B$ follow
by exchange of the two factors, and the diagonal is elementary.
For example, writing $\ell=\log(A/B)$,
\begin{align}
 \gsB_{1,1}(A,B)&=\frac{\ell}{A-B},\notag\\
 \left.\partial_u\gsB_{1+u,1+v}(A,B)\right|_{0,0}
 &=-\frac{\ell\log A+\operatorname{Li}_2(1-B/A)}{A-B},\notag\\
 \left.\partial_v\gsB_{1+u,1+v}(A,B)\right|_{0,0}
 &=-\frac{\ell\log A+\operatorname{Li}_2(1-B/A)-\ell^2/2}{A-B}.
 \label{gs:dilog-prefix}
\end{align}
Thus the finite prefix contributes explicitly reducible polylogarithms,
whereas the independent-shift information in general remains in the
two single zeta series.

\subsection{An explicit first Stieltjes coefficient}

The first spectral coefficient already combines generalized Stieltjes
constants, harmonic numbers, dilogarithms, and Hurwitz order
derivatives.  Define
\[
 D_1(A,B)=-\frac{\log(A/B)\log A+
                   \operatorname{Li}_2(1-B/A)}{A-B},
 \qquad D_1(A,A)=-\frac{1+\log A}{A},
\]
and let $\Phi^{[1]}(a,b)$ be the finite prefix
\eqref{gs:finite-prefix} with every $\gsB$ replaced by $D_1$.
These formulas are real and continuous for all $A,B>0$.

\begin{corollary}[A harmonic--Stieltjes product identity]
\label{gs:gamma1-digamma}
For $a,b>0$, choose $K,L$ satisfying
\eqref{gs:cutoff-condition} and use \eqref{gs:cd}.  Then
\begin{align}
 \int_0^1\gamma_1(a+x)\psi(b+x)\,dx
 ={}&\Phi^{[1]}(a,b)+\zeta(3)-\gamma_2(c)-\frac12\gamma_2(d)
       \notag\\
 &+\log b\,\gamma_1(c)+\zeta(2)\psi(c)
   +\left\{\zeta(2)-\frac12\log^2a\right\}\psi(d)\notag\\
 &-\sum_{j=1}^\infty\frac{(-b)^j}{j}
       \{H_j\zeta(j+1,c)+\zeta'(j+1,c)\}\notag\\
 &-\sum_{j=1}^\infty\frac{(-a)^j}{j}
          \left(\frac1j-\log a\right)\zeta(j+1,d).
 \label{gs:gamma1-digamma-formula}
\end{align}
Here $H_j=\sum_{k=1}^j1/k$ and the prime on $\zeta$ denotes its
first spectral derivative.  Both infinite sums converge geometrically.
\end{corollary}

\begin{proof}
The left side is $\partial_u\gsJ(0,0;a,b)$ under the convention
\eqref{gs:stieltjes-convention}.  The universal term contributes
$\partial_u\gsT(0,0)=\zeta(3)$, and
$\partial_u(1+u)_j|_{u=0}=j!H_j$ accounts for the harmonic number.
On the face $v=0$, the first divided difference in
\eqref{gs:completed-formula} equals
\[
 -\{\log b+\psi(1+u)+\gamma\}\gsZ(u,c)
 -\partial_u\gsZ(u,c).
\]
Its derivative at zero is
$\log b\,\gamma_1(c)+\zeta(2)\psi(c)-\gamma_2(c)$.
The derivative of the second divided difference is
$\{\zeta(2)-\tfrac12\log^2a\}\psi(d)-\tfrac12\gamma_2(d)$.
The finite prefix follows from \eqref{gs:dilog-prefix}; differentiating
the two normally convergent $\gsV$ series gives the displayed sums.
\end{proof}

\subsection{A compact logarithmic Hurwitz family}

For an integer $p\ge1$, $c>0$, $z\ge0$, set
\begin{equation}\label{gs:L-definition}
 \gsL_p(c,z)=\sum_{n=0}^\infty
       \frac{\log(1+z/(n+c))}{(n+c)^p}.
\end{equation}
This is an ordinary absolutely convergent series; it extends
holomorphically to a neighborhood of the positive real parameter
domain.  When $|z|<\Re c$, it has the geometric zeta expansion
\begin{equation}\label{gs:L-zeta}
 \gsL_p(c,z)=-\sum_{j=1}^\infty
                    \frac{(-z)^j}{j}\zeta(p+j,c).
\end{equation}
It is convenient notation for an explicit convergent series, not a
claim that this logarithmic zeta family is historically new.

Put $R_p=\partial_z\gsL_p$.  Direct partial fractions give the
terminating rules
\begin{align}
 R_1(c,z)&=\frac{\psi(c+z)-\psi(c)}z,\label{gs:R1}\\
 R_p(c,z)&=\frac{\zeta(p,c)-R_{p-1}(c,z)}z\quad(p\ge2),\label{gs:Rp}\\
 \partial_c\gsL_p(c,z)
 &=-p\gsL_{p+1}(c,z)+R_p(c,z)-\zeta(p+1,c).
 \label{gs:L-derivative}
\end{align}
All apparent singularities at $z=0$ are removable.  In particular,
\begin{equation}\label{gs:Gamma-primitive}
 \gsL_1(c,z)=\int_0^z\frac{\psi(c+y)-\psi(c)}y\,dy,
\end{equation}
with its additive constant fixed by $\gsL_1(c,0)=0$.
This identifies the precise Gamma primitive needed by the independent
shifts.  Its derivative is the divided difference of the logarithmic
derivative of $\Gamma$.

For completeness, one two-parameter member generates the full family:
\begin{equation}\label{gs:L-generator}
 \sum_{p=1}^\infty w^{p-1}\gsL_p(c,z)
 =\gsL_1(c-w,z+w)-\gsL_1(c-w,w),\qquad |w|<c,
\end{equation}
locally where all logarithms have their anchored branches.  Summing
the geometric series $\sum_{p\ge1}w^{p-1}(n+c)^{-p}$ proves this
identity.  The finite shift rule
\begin{equation}\label{gs:L-shift}
 \gsL_p(c+1,z)=\gsL_p(c,z)-c^{-p}\log(1+z/c)
\end{equation}
makes evaluation at large $z/c$ straightforward: shift $c$ until
\eqref{gs:L-zeta} converges geometrically, retaining only finitely
many elementary logarithms.

\subsection{An explicit two-shift digamma product}

Define the continuous divided difference
\[
 D(A,B)=\frac{\log(A/B)}{A-B},\qquad D(A,A)=\frac1A.
\]
For the moment suppose $a,b>0$ and $|a-b|<1$, and abbreviate
\[
 c=1+a-b,\qquad d=1+b-a.
\]

\begin{theorem}[Ordinary independent-shift digamma identity]
\label{gs:digamma}
\begin{align}
 \boxed{\begin{aligned}
 \int_0^1\psi(a+x)\psi(b+x)\,dx
 ={}&D(a,b)+\gamma_1(c)+\gamma_1(d)-2\zeta(2)\\
 &+\log b\,\psi(c)+\log a\,\psi(d)\\
 &+\gsL_1(c,b)+\gsL_1(d,a).
 \end{aligned}}
 \label{gs:digamma-formula}
\end{align}
For arbitrary $a,b>0$, choose $K\ge0,L\ge1$ so that the $c,d$ in
\eqref{gs:cd} are positive, and replace $D(a,b)$ by
\[
 \sum_{k=0}^{K-1}D(a+k,b)+\sum_{k=1}^{L-1}D(a,b+k).
\]
All other terms retain the form in \eqref{gs:digamma-formula} with
these $c,d$.
\end{theorem}

\begin{proof}
Set $u=v=0$ in \eqref{gs:completed-formula}.
The first divided difference becomes
$\gamma_1(c)-\log b\,\gamma_0(c)$, because both first derivatives
of $\gsQ\gsE$ vanish at the origin.  The second is the exchanged
expression.  Also $\gsV(0,0;c,b)=\gsL_1(c,b)$,
$\gamma_0=-\psi$, and $\gsT(0,0)=-2\zeta(2)$.
This proves the assertion under the geometric conditions.
The ordinary logarithmic series extends it to all stated positive
parameters.  Independence of the cutoffs can alternatively be checked
without analytic continuation: increasing $K$ by one adds
$D(K+a,b)$ to the prefix, while the shift formulas for $\gamma_1$,
$\psi$, and \eqref{gs:L-shift} subtract exactly that same number.
The argument for $L$ is identical.
\end{proof}

The diagonal is particularly transparent:
\begin{equation}\label{gs:diagonal-digamma}
 \int_a^{a+1}\psi(x)^2\,dx
 =\frac1a-2\gamma\log a+2\gamma_1-2\zeta(2)
                         +2\gsL_1(1,a),\qquad a>0.
\end{equation}
This specialization belongs to the established digamma-square and
logarithmic-series literature \cite{Connon}; it is included as a
normalization check, not relabelled as a new result.

\subsection{The antisymmetric first-derivative component}

Let
\begin{equation}\label{gs:W0}
 W_0(A,B)=
 \frac{2\log(A/B)-(A-B)(A^{-1}+B^{-1})}{(A-B)^2},
 \qquad W_0(A,A)=0.
\end{equation}
For $a,b>0$, $|a-b|<1$, keep $c=1+a-b,d=1+b-a$.

\begin{theorem}[Explicit antisymmetric derivative]
\label{gs:Wronskian}
The ordinary integral
\[
 W(a,b)=\int_0^1
 \{\psi(a+x)\psi'(b+x)-\psi'(a+x)\psi(b+x)\}\,dx
\]
has the exact evaluation
\begin{align}
 \boxed{\begin{aligned}
 W(a,b)={}&W_0(a,b)\\
 &+2\bigl\{\zeta'(2,d)-\zeta'(2,c)
       +\log a\,\zeta(2,d)-\log b\,\zeta(2,c)\\
 &\hspace{19mm}+\gsL_2(c,b)-\gsL_2(d,a)\bigr\}\\
 &+\frac{2\psi(c)-\psi(a+1)}b
       +\frac{\psi(b+1)-2\psi(d)}a.
 \end{aligned}}
 \label{gs:Wronskian-formula}
\end{align}
For arbitrary positive shifts, use \eqref{gs:cd}, replace $W_0(a,b)$
by its two finite-prefix sums, and replace $\psi(a+1),\psi(b+1)$
by $\psi(K+a),\psi(L+b)$ respectively.
\end{theorem}

\begin{proof}
Apply $\partial_b-\partial_a$ to \eqref{gs:digamma-formula}.
The finite divided difference gives $W_0$.
For $c=K+a-b$ one has
$(\partial_b-\partial_a)c=-2$, and the corresponding derivative
of $d$ is $2$.  Formula \eqref{gs:L-derivative} gives
\[
 (\partial_b-\partial_a)\gsL_1(c,b)
 =2\gsL_2(c,b)-R_1(c,b)+2\zeta(2,c),
\]
and the exchanged term has the opposite form.  The standard identity
\[
 \gamma_1'(z)=\zeta(2,z)+\zeta'(2,z)
\]
follows by differentiating the Hurwitz shift derivative at its spectral
pole.  It cancels the explicit $\zeta(2,c),\zeta(2,d)$ terms just
produced.  Finally use \eqref{gs:R1}, $c+b=K+a$, and $d+a=L+b$.
Every differentiation is justified by normal convergence of the
ordinary logarithmic series on compact positive parameter sets.
\end{proof}

The symmetric companion is the elementary endpoint identity
\begin{equation}\label{gs:symmetric}
 \int_0^1
 \{\psi'(a+x)\psi(b+x)+\psi(a+x)\psi'(b+x)\}\,dx
 =\frac{\psi(a)}b+\frac{\psi(b)}a+\frac1{ab}.
\end{equation}
Combining \eqref{gs:Wronskian-formula} and \eqref{gs:symmetric}
therefore evaluates each ordered first-derivative product separately.
The antisymmetric formula contains information not supplied by
integration by parts alone.

\subsection{Closure under every ordinary argument derivative}

Write $q=r+t$.  In differentiating \eqref{gs:digamma-formula}, the
only non-elementary series that must remain after $q$ derivatives are
\begin{equation}\label{gs:highest-L}
 q!\{(-1)^r\gsL_{q+1}(c,b)+(-1)^t\gsL_{q+1}(d,a)\}.
\end{equation}
All other terms reduce finitely to polygamma values, rational
functions, logarithms, and Hurwitz-zeta order derivatives.  For
$q>0$, even the two generalized Stieltjes constants disappear:
\begin{equation}\label{gs:gamma1-derivatives}
 \gamma_1^{(q)}(z)=(-1)^{q+1}q!
        \{\zeta'(q+1,z)+H_q\zeta(q+1,z)\}.
\end{equation}
To prove \eqref{gs:highest-L}, note that the highest remaining
$\gsL$ term in a derivative with respect to $a$ of $\gsL_p(c,b)$
is $-p\gsL_{p+1}$, whereas a derivative with respect to $b$ has
highest term $p\gsL_{p+1}$.  The remaining $R_p$ terms terminate
by \eqref{gs:R1}--\eqref{gs:Rp}.  Induction gives the displayed
coefficient.  The other term is exchanged.  Equation
\eqref{gs:gamma1-derivatives} follows by taking the coefficient of
the linear spectral term in
$\partial_z^q\zeta(1+u,z)=(-1)^q(1+u)_q\zeta(q+1+u,z)$.
Thus the rules are a terminating reduction algorithm, rather than an
instruction to leave arbitrary derivatives unevaluated.

For example, a completely explicit trigamma-square identity is
\begin{align}
 \int_a^{a+1}\psi'(x)^2\,dx
 ={}&\frac1{3a^3}+4\zeta'(3)+\frac{4\zeta(2)}a
        +4\zeta(3)\log a\notag\\
 &-\frac{2\{\psi(1+a)+\gamma\}}{a^2}
        -4\gsL_3(1,a),\qquad a>0.
 \label{gs:trigamma-square}
\end{align}
It is obtained by differentiating the two independent shifts first
and only then setting $a=b$.  Merely differentiating
\eqref{gs:diagonal-digamma} twice would yield the sum of several
ordered derivative products, not this individual square.

\subsection{Diagonal limits and the inherited coincident finite parts}

There are two different diagonal operations.  Setting $a=b>0$
is an ordinary holomorphic specialization of
\eqref{gs:completed-formula}.  Sending the common shift to zero is
a boundary operation and requires removal of the local divergences.
For fixed indices $m,n,r,t$, write
\[
 P(x)=\gamma_m^{(r)}(x)\gamma_n^{(t)}(x).
\]
The exact recurrence
$\gamma_m(x)=x^{-1}\log^m x+\gamma_m(1+x)$ supplies its finite
singular expansion
\[
 P(x)=\sum_{\alpha\ge1,h\ge0}C_{\alpha,h}x^{-\alpha}\log^h x
           +O(1+|\log x|^M)
\]
near zero, with finitely many terms in the displayed singular sum.
The coefficients are explicit products of the Laurent singular
polynomials and Taylor derivatives at one.  For each monomial let
$A_{\alpha,h}(x)$ be an antiderivative with zero constant term in
its power-log expansion.  Then
\begin{equation}\label{gs:boundary-limit}
 \lim_{a\downarrow0}\left\{
 \int_0^1P(a+x)\,dx
       +\sum_{\alpha,h}C_{\alpha,h}A_{\alpha,h}(a)\right\}
 =\operatorname{FP}\int_0^1P(x)\,dx.
\end{equation}
Indeed, the shifted integral is $\int_a^{1+a}P$, and its difference
from $\int_a^1P$ tends to zero.  After integration, the singular
monomials contribute $-C_{\alpha,h}A_{\alpha,h}(a)$; the remaining
integral converges by the displayed local bound.  This proves the
boundary compatibility for every fixed order, in the original
coordinate $x$, without exchanging an unjustified spectral limit
with a finite-part operation.

For example, \eqref{gs:diagonal-digamma} and
\eqref{gs:trigamma-square} give the exact checks
\begin{align}
 \operatorname{FP}\int_0^1\psi(x)^2\,dx
     &=2\gamma_1-2\zeta(2),\notag\\
 \operatorname{FP}\int_0^1\psi(x)\psi'(x)\,dx
     &=\zeta(2),\notag\\
 \operatorname{FP}\int_0^1\psi'(x)^2\,dx
     &=2\zeta(3)+4\zeta'(3).
 \label{gs:boundary-checks}
\end{align}
These agree with the inherited coincident product generator.  They
are compatibility checks, not newly claimed coincident evaluations.
Unequal-rate collisions $a,b\to0$ are not covered by a claim that
one can omit the local subtractions; their path-dependent corrections
remain a separate question.

\subsection{Source reconciliation and further targets}

There is a concrete status correction to the incoming research list.
Question R7 of \emph{Tail Kernels and Boundary Products} asks for all
nonpositive-order continuations of the alternating square law.
The cutoff-jet section of the independently archived
\emph{Bilinear Tails and Determinants} already proves the requisite
normal subtractions, all cutoff contacts, their spectral jets, the
order-zero $1/4$ discrepancy, and the change under a translated cutoff
\cite{RepoBilinear}.  R7 should therefore be marked resolved by
that report, with its stated coordinate convention.  Its Gamma-square
specializations should not be counted again as new results here.

The completed generator resolves R8 in the explicit convergent function
class described above.  It does not claim that arbitrary independent
shifts reduce to finitely many ordinary zeta constants.  Three sharper
questions remain:
\begin{enumerate}
\item Determine special rational pairs $(a,b)$ for which the two
logarithmic Hurwitz series in \eqref{gs:Wronskian-formula} reduce
further to a finite collection of classical polylogarithms and Gamma
values.  The coincident and complementary-shift cases should be
separated from generic rational pairs.
\item Derive functional equations in the shift difference $a-b$
that relate finite descriptions of the prefix by polylogarithms
to the two geometric zeta series.  The exact $K,L$ invariance gives
necessary tests for any proposed relation.
\item Compute the local correction for $a=\lambda\epsilon$,
$b=\mu\epsilon$ with fixed positive $\lambda,\mu$ and compare its
constant term with the coincident finite part.  The generator makes
this a defined two-parameter collision problem, but the diagonal
compatibility theorem alone does not settle it.
\end{enumerate}
